%%Cor_deltaFF
Selon la valeur du flag :\\
1- si le flag est �gal � 0 (mesures ok) alors : \\
\hspace*{2cm} $flag\_Cdf$ = 0 et $ffbas\_Cdf$ = $ffbas$ etc...
\\
\\
2- si le flag est positif ($ffhaut$ trop faible par rapport � $ffbas$) alors : \\
\\
\indent Si le point est isol� (flag avant et apr�s �gaux � 0) , on recalcule $ffhaut\_Cdf$ � partir des valeurs interpol�s selon la formule suivante : \\
\hspace*{2cm} $coeff$ = (( $ffhaut_{i-1}$ / $ffbas_{i-1}$ ) + ( $ffhaut_{i+1}$/$ffbas_{i+1}$ )) / 2.  \\
\hspace*{2cm} $ffhaut\_Cdf$ = ( 1. / $coeff$ ) * $ffbas$  \\
\hspace*{2cm} $ffbas\_Cdf$ = $ffbas$  \\
\hspace*{2cm} $flag\_Cdf$ = 1 \\
\\
\indent Si le point n'est pas isol� (flag avant et/ou apr�s diff�rent de 0) , on recalcule $ffhaut\_Cdf$ par la m�thode empirique qui a servi � d�tecter les points de mesures incoh�rents entr'eux : \\
\hspace*{2cm} $ffhaut\_Cdf$ = ( 1.1647 * $ffbas$ ) + 0.1  \\
\hspace*{2cm} $ffbas\_Cdf$ = $ffbas$  \\
\hspace*{2cm} $flag\_Cdf$ = 2 \\
\\
\\
3- si le flag est n�gatif ($ffbas$ trop faible par rapport � $ffhaut$) alors : \\
\\
\indent Si le point est isol� (flag avant et apr�s �gaux � 0) , on recalcule $ffbas\_Cdf$ � partir des valeurs interpol�s selon la formule suivante : \\
\hspace*{2cm} $coeff$ = (( $ffhaut_{i-1}$ / $ffbas_{i-1}$ ) + ( $ffhaut_{i+1}$/$ffbas_{i+1}$ )) / 2.  \\
\hspace*{2cm} $ffbas\_Cdf$ = $coeff$ * $ffbas$  \\
\hspace*{2cm} $ffhaut\_Cdf$ = $ffhaut$  \\
\hspace*{2cm} $flag\_Cdf$ = 1 \\
\\
\indent Si le point n'est pas isol� (flag avant et/ou apr�s diff�rent de 0) , on recalcule $ffbas\_Cdf$ par la m�thode empirique qui a servi � d�tecter les points de mesures incoh�rents entr'eux : \\
\hspace*{2cm} $ffbas\_Cdf$ = ( 0.8586 * $ffhaut$ ) - 0.08586  \\
\hspace*{2cm} $ffhaut\_Cdf$ = $ffhaut$  \\
\hspace*{2cm} $flag\_Cdf$ = 2 \\
\\
\\
RAPPEL : Formules des droites empiriques d�limitant la zone de coh�rence utilis�es pour la "d�tection" des anomalies (op�rateur Valid\_empir\_deltaFF) :\\
\hspace*{1cm} Si $ffbas$ < 4 m/s :\\
\hspace*{2cm} $ffhaut$ doit �tre inf�rieur � $ffhaut_{maxi1}$ avec \\ 
$ffhaut_{maxi1}$ = 1.225 * $ffbas$ + 0.7\\
\\
\hspace*{1cm} Sinon :\\
\hspace*{2cm} $ffhaut$ doit �tre inf�rieur � $ffhaut_{maxi2}$ avec \\ 
$ffhaut_{maxi2}$ = 1.6 * $ffbas$ - 0.8\\
\\
\hspace*{1cm} Dans tous les cas (quelle que soit la valeur de $ffbas$):\\
\hspace*{2cm} $ffhaut$ doit �tre sup�rieur � $ffhaut_{min}$ avec \\ 
$ffhaut_{min}$ = 1.08 * $ffbas$ - 0.5\\
